\usepackage{fontspec}
\setmonofont{Courier New}                          % моноширинный шрифт
\newfontfamily\cyrillicfonttt{Courier New}         % моноширинный шрифт для кириллицы
\defaultfontfeatures{Ligatures=TeX}                % стандартные лигатуры TeX, замены нескольких дефисов на тире и т. п. Настройки моноширинного шрифта должны идти до этой строки, чтобы при врезках кода программ в коде не применялись лигатуры и замены дефисов
\setmainfont{Times New Roman}                      % Шрифт с засечками
\newfontfamily\cyrillicfont{Times New Roman}       % Шрифт с засечками для кириллицы
\setsansfont{Arial}                                % Шрифт без засечек
\newfontfamily\cyrillicfontsf{Arial}               % Шрифт без засечек для кириллицы

% \usepackage{amsfonts,amssymb} убрано - несовместимо с unicode-math

\usepackage[math-style=TeX]{unicode-math}
%\usepackage{lua-unicode-math}
\setmathfont{STIX Two Math}
%\setmathfont{XITS Math}

%%% Правила русской математической типографии %%%
% Наклонные неравенства (slanted inequalities)
\AtBeginDocument{%
  \renewcommand{\le}{\leqslant}%
  \renewcommand{\leq}{\leqslant}%
  \renewcommand{\ge}{\geqslant}%
  \renewcommand{\geq}{\geqslant}%
}

% var-версии греческих букв по умолчанию (нужно только при использовании unicode-math)
\AtBeginDocument{%
  \renewcommand{\epsilon}{\varepsilon}%
  \renewcommand{\phi}{\varphi}%
  \renewcommand{\kappa}{\varkappa}%
  \renewcommand{\theta}{\vartheta}%
}

% Шрифты должны загружаться ДО polyglossia
\usepackage{polyglossia}
\setmainlanguage[babelshorthands=true]{russian}
\setotherlanguage{english}
